\subsection{Properties true almost everywhere}

Let X be a locally compact space, \(\mu\) a measure on X. If \(P\{x\}\) is a property, the property \(\langle P\{x\}\) almost everywhere (with respect to \(\mu\) ) is by definition equivalent to the property \(\langle\) the set of x such that \((x \in X \text{ and not } P\{x\})\) is \(\mu\) -negligible.

THEOREM 1. --- In order that a numerical function (finite or not) \(f \geqslant 0\) defined on X be negligible, it is necessary and sufficient that \(f(x) = 0\) almost everywhere.

The condition is necessary. For, suppose that \(f\) is negligible, and let \(N\) be the set of \(x \in X\) such that \(f(x) \neq 0\) ; then \(\varphi_N \leqslant \sup_n (nf)\) , therefore \(\varphi_N\) is negligible (No.~1, Props. 1 and 2).

The condition is sufficient. Suppose that the set N of points where \(f(x) \neq 0\) is negligible; then \(f \leqslant \sup_{n} n\varphi_{N}\) , therefore f is negligible (No.~1, Props. 1 and 2).

PROPOSITION 6. --- If \(f\) and \(g\) are two functions \(\geqslant 0\) (finite or not) defined on X such that \(f(x) = g(x)\) almost everywhere, then \(|\mu|^*(f) = |\mu|^*(g)\) .

Let N be the negligible set of points \(x \in X\) such that \(f(x) \neq g(x)\) . The functions \(\inf(f, g)\) and \(\sup(f, g)\) being equal except at the points of N, it suffices to prove the proposition assuming \(f \leqslant g\) . Let h be the function equal to \(+\infty\) at the points of N, and to 0 on CN; then \(f \leqslant g \leqslant f + h\) , thus

\[
| \mu | ^ {*} (f) \leqslant | \mu | ^ {*} (g) \leqslant | \mu | ^ {*} (f + h) \leqslant | \mu | ^ {*} (f) + | \mu | ^ {*} (h) = | \mu | ^ {*} (f)
\]

(since \(h\) is negligible), whence the proposition.

PROPOSITION 7. --- If \(f\) is a function \(\geqslant 0\) defined on \(X\) such that \(|\mu|^*(f) < +\infty\) , then \(f(x)\) is finite almost everywhere.

For, let N be the set of points \(x \in X\) such that \(f(x) = +\infty\) ; for every integer n, \(n\varphi_{N} \leqslant f\) , whence \(n|\mu|^{*}(\varphi_{N}) \leqslant |\mu|^{*}(f)\) ; since n is arbitrarily large, \(|\mu|^{*}(N) = 0\) .

However, even if X is compact, a function \(f \geqslant 0\) defined on X and everywhere finite can have infinite upper integral, as is shown by the example X = {[}0, 1{]}, \(f(x) = 1/x\) for x \textgreater{} 0 and \(f(0) = 0\) , \(\mu\) being Lebesgue measure on X.

